\section{Gluing and links}

\begin{statement}{C.1}{Proved}{The categorical gluing datum is a bimodule}
Let $A\subset X$ be a full sieve in a small category, let $B=X\setminus A$ be its full complementary cosieve, and write $j:A\hookrightarrow X$ and $i:B\hookrightarrow X$. The datum of $X$ with these specified inclusions is equivalent to the datum of $A,B$ and
\[
 M:A^{\op}\times B\longrightarrow\Set,
 \qquad M(a,b)=\Hom_X(a,b).
\]
Its category of crossing arrows $E$ has objects $(a,b,m)$, $m\in M(a,b)$, and arrows $(u,v):(a,b,m)\to(a',b',m')$ satisfying $u:a\to a'$, $v:b\to b'$ and $vm=m'u$. Its source and target are functors $s:E\to A$ and $t:E\to B$.
\end{statement}
\begin{proof}
There is no arrow from $B$ to $A$, since $A$ is a sieve. All remaining hom sets are those of $A$, those of $B$, and the values of $M$. The functoriality of $M$ gives precomposition by arrows of $A$ and postcomposition by arrows of $B$. Its interchange identity is precisely associativity for a composite containing one crossing arrow. A composite can contain at most one crossing arrow, so there are no other associativity conditions. Conversely these operations recover the stated functor. On morphisms one uses functors on $A,B$ and a natural transformation of their crossing bimodules. Thus the reconstruction is an equivalence of categories of specified gluing data, with the evident canonical isomorphism back to the original $X$. This verifies XII ms.~59--60; the printed claim of an isomorphism of categories requires chosen object sets, as the editorial correction already notes. No localization or choice of paths occurs here.
\end{proof}

\begin{statement}{C.2}{Known}{The crossing category is the homotopy link, with the correct complementary open subset}
Form $K=N(X^{\op})\to\Delta[1]$ by assigning $B$ to $0$ and $A$ to $1$. Give $Z=|K|$ the last-vertex stratification over $\{0<1\}$ with its upper Alexandrov topology. Then its closed stratum is $Z_0=B(B^{\op})$ and its open stratum is $Z_1=Z\setminus Z_0$. There is a natural weak equivalence
\[
 B(E^{\op})\longrightarrow
 \operatorname{holink}(Z_0,Z_1),
\]
where a homotopy-link path starts in $Z_0$ and lies in $Z_1$ at every positive time. Under the canonical opposite-nerve homeomorphisms, $Z$ is $BX$, $Z_0$ is $BB$, and $Z_1$ deformation retracts onto $BA$. The comparison is compatible with both endpoint maps, followed at the open end by this deformation retraction.
\end{statement}
\begin{proof}
An $n$-simplex of $\operatorname{Map}_{\Delta[1]}(\Delta[1],K)$ is a functor $[n]\times[1]\to X^{\op}$ which has its first row in $B^{\op}$ and its second in $A^{\op}$. Reversing its arrows identifies it with an $n$-simplex of $N(E^{\op})$. This is an equality of simplicial sets, including faces and degeneracies. The realization map to topological homotopy links is the weak equivalence of \citet[Theorem~6.1]{CDDouteauWaas2022}, applied to this particular stratified simplicial set. That theorem applies to arbitrary stratified simplicial sets; no EI, finiteness, or Quinn-fibration condition has been added implicitly.

For the claimed retraction, represent a point by barycentric coordinates on a simplex of $N(X^{\op})$. A point belongs to $Z_1$ exactly when the total weight $w$ of its $A$ vertices is positive. Multiply every $B$ coordinate by $1-r$ and every $A$ coordinate by $(1-r)+r/w$, for $0\leq r\leq1$. The sum remains one, the coordinates remain nonnegative, and the formula commutes with faces and degeneracies. It defines a strong deformation retraction of $Z_1$ onto $B(A^{\op})$. The realization comparison in the cited theorem is obtained by evaluating the prism $|\Delta[n]\times\Delta[1]|$; its endpoints are therefore the source and target evaluations, not an unspecified abstract weak equivalence. Opposite-nerve homeomorphisms reverse the order of barycentric coordinates and identify these maps with $Bs$ and $Bt$.

For uncountable categories $Z$ need not be metrizable, so it need not be a homotopically stratified set in Quinn's original sense. Quinn's tameness and homotopy-link fibration hypotheses remain genuine additional hypotheses \citep[\S\S2.1,3.1]{CDQuinn1988}. The theorem just used is the simplicial/stratified-space form of the comparison and avoids that unnecessary restriction. In particular $BA$ and $BB$ are not being asserted to be literal complementary subsets of $BX$: both are subcomplexes, and the complementary open set is $BX\setminus BB$.
\end{proof}

\begin{statement}{C.3}{Known}{The specified categorical and homotopy-link gluing maps agree}
For the data of C.1, the comparison determined by the natural transformation $js\Rightarrow it$ is a weak homotopy equivalence
\[
 \hocolim(BA\xleftarrow{Bs}BE\xrightarrow{Bt}BB)\longrightarrow BX.
\]
Via C.2 it is the homotopy-link double-mapping-cylinder comparison. This includes empty $A$, empty $B$, empty $E$, and arbitrarily many connected components.
\end{statement}
\begin{proof}[Verification of the supplied proof and comparison]
We use the supplied POLES G.19 proof rather than give a third proof. Its category $P$ has the three fibres $A,E,B$, arrows $a'\to(a\to b)$ indexed by $a'\to a$, and arrows $b'\to(a\to b)$ indexed by $b'\to b$. The displayed functor $P\to X$ sends a crossing-arrow object to its target. For $x\in A$, its comma category is $A/x$. For $x\in B$, the endofunctor printed there sends $a\to x$ to the corresponding crossing arrow, sends $(a\to b\to x)$ to its composite, and sends an object of $B/x$ to $\id_x$. The two printed natural transformations are well typed: on a crossing object their component is the square $(\id_a,b\to x)$, and on the other fibres it is the indicated structural arrow. The squares commute by the definition of $E$; the constant functor involved takes its value at the $B$ object $(x,\id_x)$. Thus every comma category is contracted exactly as claimed. The fibred Grothendieck-construction comparison identifies $BP$ with the stated homotopy colimit and identifies its map to $BX$ with the source-to-target homotopy. These checks validate the application of Theorem A in POLES G.19 \citep{CDPOLES2026}.

The same source-to-target prism occurs in C.2. After reversing its interval coordinate to start at $BB$, and applying the deformation retraction at its open endpoint, it is the holink evaluation prism. Hence the comparison is the same coherent cocone. This derives the weak homotopy-link gluing assertion for these nerves from the audited categorical proof. It does not assert Quinn's stronger complement-preserving homotopy equivalence without his tameness hypotheses. If one of the pieces is empty, the corresponding crossing category is empty and the homotopy colimit reduces to the remaining piece. If $E$ alone is empty, $X=A\amalg B$ and it reduces to a disjoint union.
\end{proof}

\begin{statement}{C.4}{Refuted}{Unrestricted paths cannot replace the crossing link}
Replacing the crossing category in C.3 by all paths between its endpoint subspaces changes the answer, already for a category with two objects and four arrows.
\end{statement}
\begin{proof}
Let $X$ have objects $a,b$, identity arrows and exactly two further arrows $r,s:a\to b$. Composition with identities is the only composition. Take $A=\{a\}$ and $B=\{b\}$. Then $E$ is the discrete two-object category and $BX$ is the circle obtained from two edges with their endpoints identified in pairs. Consequently the unrestricted path space from $a$ to $b$ is homotopy equivalent to $\Omega S^1$ and has components indexed by $\mathbb Z$, each contractible. In contrast $BE$ has exactly two components. The double mapping cylinder of two points under $\mathbb Z$ has first integral homology freely generated by countably many differences of edges, whereas the double mapping cylinder under $\{r,s\}$ has first homology $\mathbb Z$. Thus even the resulting unstratified gluing type changes. The category and its nerve boundary matrix are checked in \texttt{code/CD\_checks.py}; its output records Betti numbers $(1,1)$. The distinction is between paths constrained to leave the closed stratum immediately and paths allowed to return to it and acquire arbitrary winding.
\end{proof}

\begin{statement}{C.5}{Proved}{Finite filtrations require the actual crossing categories and their flag data}
Let $\varnothing=A_0\subset A_1\subset\cdots\subset A_r=X$ be a finite filtration by full sieves, and put $B_k=A_k\setminus A_{k-1}$. At stage $k$ let $E_k$ be the category of arrows in $A_k$ from $A_{k-1}$ to $B_k$. Then
\[
 BA_k\simeq\hocolim(BA_{k-1}\leftarrow BE_k\to BB_k)
\]
by the specified comparisons of C.3. All generalized links are given by the nerves of the categories of flagged strings of arrows: for $i_0<\cdots<i_m$, their objects are $b_0\to\cdots\to b_m$ with $b_j\in B_{i_j}$ and their morphisms are commuting ladders. These categories, their restriction functors, and their composition functors retain the iterated gluing information.
\end{statement}
\begin{proof}
Each $A_{k-1}$ is a sieve in $A_k$, so C.3 applies successively. This is an induction using actual categories and actual functors. There is no appeal to an isomorphism of pro-objects or to a limit of homotopy categories. The level function $X\to\{1<\cdots<r\}$ is a functor; its opposite stratifies $N(X^{\op})$ over the oppositely ordered level set. Applying the simplicial calculation of C.2 to the reversed flag instead of $[1]$ identifies its simplicial mapping space with the nerve of the opposite of the specified ladder category. Douteau--Waas Theorem~6.1 again supplies its geometric interpretation. Restriction of strings supplies the face maps, and ordinary composition in $X$ supplies the remaining flag maps; all their identities hold before realization.

One must not replace these data by pairwise homotopy types without their coherence. Even the boundary of a triangle and the triangle, stratified by their ordered vertices, have equivalent individual strata and pairwise holinks but different three-stratum links: the full-flag holink is empty for the boundary and nonempty for the triangle \citep[Example~2.30]{CDDouteauWaas2022}. Their underlying spaces also have different first homology. The accompanying program checks the boundary matrices. When every layer is a groupoid, the ladder categories are groupoids and their endpoint maps are isofibrations (lift an endpoint isomorphism by whiskering); their strict Segal fibre products then are homotopy fibre products. This is the spatial d\'ecollage situation. For arbitrary non-groupoidal layers one retains the flagged categories themselves; groupoidification has not been claimed to preserve arbitrary strict fibre products.
\end{proof}

\begin{statement}{C.6}{Proved}{Enhanced coefficients glue by a lax pullback}
Let $\cC$ be a presentable $\infty$-category, and put $\mathcal F_C=\Fun(C^{\op},\cC)$ for a small category $C$. Define
\[
 T=i^*j_*:\mathcal F_A\longrightarrow\mathcal F_B.
\]
Then $\mathcal F_X$ is equivalent to the $\infty$-category of triples $(F_A,F_B,\theta:F_B\to TF_A)$. There is a canonical equivalence $T\simeq t_*s^*$, so the gluing arrow may equivalently be specified as a homotopy-coherent natural transformation $t^*F_B\to s^*F_A$ in $\mathcal F_E$. At $b\in B$,
\[
 (TF_A)(b)\simeq\limop_{(A/b)^{\op}}F_A.
\]
The limit is an $\infty$-categorical limit; neither it nor the lax pullback may be replaced by its ordinary counterpart in homotopy categories.
\end{statement}
\begin{proof}
An object $F\in\mathcal F_X$ has the unit $F\to j_*j^*F$, whose restriction to $B$ gives $\theta$. Conversely define, from the given coherent arrow, an object of $\mathcal F_X$ by
\[
 F=j_*F_A\mathop{\times}_{i_*TF_A}i_*F_B.
\]
The map from the first factor is the unit of $i^*\dashv i_*$ and the map from the second is $i_*\theta$. Restriction to $A$ gives $F_A$, since $i_*$ has terminal values there. Restriction to $B$ gives $TF_A\times_{TF_A}F_B\simeq F_B$. Applied to an original $F$, its two units therefore give a comparison to this pullback which is an equivalence at every object, hence an equivalence of functors. Applied to a triple, the same restrictions recover its specified arrow. The constructions use functorial limits and units, so these are inverse equivalences of $\infty$-categories, including mapping spaces.

To identify $T$, consider $r:E/b\to A/b$, $(a\xrightarrow{m}b',\beta:b'\to b)\mapsto(a,\beta m)$. It has a fully faithful right adjoint $k:A/b\to E/b$, $(a,n:a\to b)\mapsto(n,\id_b)$. The adjunction is the equality between a square with target component $\beta$ and its source component $a\to a'$. Thus $k$ is homotopy final, $k^{\op}$ is homotopy initial, and the pointwise right-Kan formula identifies $t_*s^*F_A$ with the displayed limit over $(A/b)^{\op}$. Adjunction $t^*\dashv t_*$ now identifies the two spaces of gluing arrows, naturally in both coefficient objects. The ordinary set-valued case is the gluing of topoi reproduced at XII ms.~62--63 from SGA~4, Expos\'e~IV, Theorem~9.5.4 and Exercise~9.5.10 \citep{CDSGA4}. The argument above proves the enhanced version directly, rather than inferring it from ordinary sheaf gluing.
\end{proof}

\begin{statement}{C.7}{Refuted}{The ordinary comma category of homotopy categories loses gluing morphisms}
Even for $X=[1]$ and space-valued coefficients, the homotopy category of C.6 is not reconstructed by replacing each $\infty$-category in its lax pullback by its homotopy category and taking an ordinary comma category.
\end{statement}
\begin{proof}
View two objects of $\Fun([1]^{\op},\Spaces)$ as the arrows $f:S^0\to *$ and $g:*\to S^1$. Their mapping space is
\[
 \Map(S^0,*)\mathop{\times}^{h}_{\Map(S^0,S^1)}\Map(*,S^1)
 \simeq *\mathop{\times}^{h}_{S^1\times S^1}S^1,
\]
where the last map is the diagonal. Its connected components are the cokernel of $\mathbb Z\xrightarrow{(1,1)}\mathbb Z^2$, hence form $\mathbb Z$. In the ordinary arrow category of $\Ho\Spaces$, both vertical maps in such a square have a unique homotopy class and their composites agree, so there is exactly one morphism from $f$ to $g$. Thus the comparison is not faithful. The integer boundary-rank check is included in the program. For a bare derivator, the Kan extensions, support subderivators, and the pointed recollement proved in Part~D remain meaningful; the false ordinary comma formula is not a substitute for an enhancement.
\end{proof}

\begin{statement}{C.8}{Proved}{Locally constant enhanced coefficients glue by an ordinary $\infty$-categorical pullback}
For any complete and cocomplete $\infty$-category $\cC$, the locally constant part is
\[
 \Fun(BX,\cC)\simeq
 \Fun(BA,\cC)\mathop{\times}_{\Fun(BE,\cC)}\Fun(BB,\cC),
\]
where the pullback is in $\Cat_\infty$ and therefore uses an equivalence between the two restricted local systems. The same assertion holds for every stage of C.5.
\end{statement}
\begin{proof}
Groupoidification identifies locally constant functors on $X^{\op}$ with functors on $BX$. C.3 is a pushout in spaces, and its image is a pushout in $\Cat_\infty$: the inclusion of spaces is left adjoint to the maximal-subgroupoid functor and hence preserves colimits. Functor categories turn a colimit in their first variable into a limit, giving the stated equivalence. More explicitly, a functor out of the double mapping cylinder consists of two functors and a natural equivalence along $BE$; all their higher compatibility is part of the mapping-space pullback. The same argument applied at each stage proves the finite-filtration assertion. Taking $\Ho$ after forming this pullback gives the locally constant homotopy category, but does not turn it into an ordinary categorical pullback. The abstract-derivator invariance theorems of the supplied GALOIS B.3 and B.5 do not change that distinction.
\end{proof}

\begin{statement}{C.9}{Proved}{The exact weight test for a homotopically valid strict enriched gluing}
For a sieve $U\hookrightarrow X$ let $H_t(u)=\Hom_X(u,t)$, $t\in X\setminus U$. Regard ordinary categories as simplicial categories with discrete hom objects. Cobase change along this sieve preserves all Dwyer--Kan equivalences of simplicial categories under $U$ if and only if each degreewise coend
\[
 H_t\otimes_U-:\Fun(U,\mathbf{sSet})\longrightarrow\mathbf{sSet}
\]
preserves objectwise weak homotopy equivalences. This condition, discrete flatness, is stronger than the homotopical cofibration condition for classifying spaces.
\end{statement}
\begin{proof}
In the pushout against $p:U\to C$, the new hom object from $c\in C$ to $t$ is precisely $H_t\otimes_U C(c,p(-))$; all other hom objects are unchanged or empty. This follows by writing a crossing composite as one arrow in $C$, one arrow in $X$ crossing the sieve, and imposing the usual balanced relations. Thus preservation by all weights gives full faithfulness up to weak equivalence after pushout. Essential surjectivity follows from that of the original Dwyer--Kan equivalence, because the complementary objects are unchanged.

Conversely, for any $L:U\to\mathbf{sSet}$ adjoin a new source $c$ to $U$, with hom objects $C_L(c,u)=L(u)$ and no nonidentity arrows into $c$. An objectwise weak equivalence $L\to L'$ gives a Dwyer--Kan equivalence $C_L\to C_{L'}$. After pushout its hom map from $c$ to $t$ is the weight map under examination. Preservation of every Dwyer--Kan equivalence therefore forces the asserted condition. This independently checks the proof of TOTASPH F.14 \citep{CDTOTASPH2026}.

For the claimed implication to homotopical cofibrations we use the generalized version of \citet[Theorem~1.5, \S4, p.~2180]{CDHORR2024}: for a discretely flat inclusion the comparison $NX\amalg_{NU}NC\to N(X\amalg_U C)$ is a categorical equivalence. Since $NU\to NX$ is a monomorphism, its source computes the homotopy pushout of spaces. Comparing this construction for a weak equivalence $C\to C'$ under $U$ and using homotopy-pushout invariance proves the implication. C.11 proves that the converse fails. In particular a homotopy colimit in C.3 has not been exchanged for a colimit of weights without a hypothesis.
\end{proof}

\begin{statement}{C.10}{Proved}{A universally harmless attachment can have a disconnected individual crossing category}
There is a three-object, seven-arrow category $X$ and a one-object sieve $U$ such that $U\to X$ remains a weak equivalence after every pushout, while $U/t$ has two components for a complementary object $t$.
\end{statement}
\begin{proof}
Take objects $a,t,s$, nonidentity arrows $\alpha,\beta:a\to t$, $\gamma:t\to s$, and $\delta:a\to s$, with $\gamma\alpha=\gamma\beta=\delta$. Set $U=\{a\}$. These equations and identities specify all compositions. For a functor $U\to C$ selecting $c$, the pushout $P$ has the old category $C$, the two new objects $t,s$, hom sets $P(d,t)=C(d,c)\times\{\alpha,\beta\}$, $P(d,s)=C(d,c)$, and $P(t,s)=\{\gamma\}$. There are no reverse crossing arrows. The retraction $r:P\to C$ sends both new objects to $c$ and forgets the two-arrow label. The endofunctor $F:P\to P$ fixes $C$ and sends $t,s$ to $s$. There are natural transformations $\id_P\Rightarrow F\Leftarrow jr$, with components $\gamma,\id_s$ for the first and the arrow corresponding to $\id_c$ for the second; both are identities on $C$. Naturality at $\alpha,\beta$ is exactly their equality after composition with $\gamma$. Thus $Bj$ is a homotopy equivalence for every $C$ and every attaching functor. Yet $U/t$ is discrete with two objects. The program checks all compositions, associativity, the sieve condition, both transformations on the original category, and the crossing components. This verifies the universal argument in TOTASPH F.18, rather than inferring universality from finite testing. The source $U$ is discrete, so its weights are coproducts of copies and satisfy C.9.
\end{proof}

\begin{statement}{C.11}{Proved}{A homotopical cofibration need not be discretely flat}
There is a four-object, eleven-arrow full sieve $U\hookrightarrow X$ which is a couniversal weak equivalence but fails the weight test of C.9.
\end{statement}
\begin{proof}
Let $U=(0\xrightarrow{a}1)$. Adjoin $t$ and arrows $\beta,\gamma:1\to t$ and $c:0\to t$, with $\beta a=\gamma a=c$, and then adjoin a terminal object $s$. There are no other arrows except identities and the unique arrows into $s$. Send $t,s$ to $1$ to obtain a retraction $r:X\to U$. The endofunctor fixing $U$ and sending $t,s$ to $s$ gives $\id_X\Rightarrow F\Leftarrow ir$, relative to $U$. On pushing out along any $U\to C$, these functors and transformations extend by the identity on $C$: their component equations are unchanged balanced-composite equations. Hence each pushout inclusion is a categorical homotopy equivalence. Consequently cobase change preserves weak equivalences under $U$ by two-out-of-three.

For non-flatness take a simplicial circle $K$ and its cone $CK$. On $U$ define $L=(K\hookrightarrow CK)$ and $L'=(K\to *)$, with comparison identity at $0$ and the contraction at $1$. This is an objectwise weak equivalence. The crossing weight $H_t$ is singleton at $0$ and has two elements at $1$, both restricting to the singleton. Directly from the coend relations,
\[
 H_t\otimes_U L=CK\amalg_K CK,
 \qquad H_t\otimes_U L'=*.
\]
The first is the suspension of the circle, with integral $H_2=\mathbb Z$; its map to the second is not a weak equivalence. The program checks the finite category, its relative transformations, and a simplicial two-cone model with Betti numbers $(1,0,1)$. This completes the proposed TOTASPH F.21 test. It also explains why ``the strict gluing has the right classifying space'' and ``all crossing weights compute their derived values strictly'' are different assertions.
\end{proof}
